%%%%%%%%%%%%%%%%%%%%%%%%%%%%%%%%%%%%%%%%%%%%%%%%%%%%%%%%%%%%%%%%%%%%%%
% Lớp 8 | Chương I. Đa thức | Phép nhân đa thức
% Dạng 8D1?4-3: Rút gọn, chứng minh biểu thức không phụ thuộc biến
%%%%%%%%%%%%%%%%%%%%%%%%%%%%%%%%%%%%%%%%%%%%%%%%%%%%%%%%%%%%%%%%%%%%%%


% nguồn: De_AI_Toan8_cuoi_hoc_ki_i_1 (nhập 21/11/2026)
\begin{bt}%[8D1H4-3]%[ma:fc96431b]
	Chứng minh rằng giá trị của biểu thức $D = (x - 3)(x + 5) - (x - 1)(x + 3)$ không phụ thuộc vào giá trị của biến $x$.
	\loigiai{
		Ta có $D = (x^2 + 5x - 3x - 15) - (x^2 + 3x - x - 3)$ \\
$= (x^2 + 2x - 15) - (x^2 + 2x - 3) = x^2 + 2x - 15 - x^2 - 2x + 3 = -12$. \\
Vì $D = -12$ không chứa biến $x$ nên giá trị của $D$ không phụ thuộc vào biến $x$.
	}
\end{bt}

% nguồn: De_AI_Toan8_cuoi_hoc_ki_i_1 (nhập 21/11/2026)
\begin{bt}%[8D1H4-3]%[ma:53aa2428]
	Chứng minh rằng giá trị của biểu thức $E = (2x - 1)(3x + 2) - 6x(x + 1) + 5x$ không phụ thuộc vào giá trị của biến $x$.
	\loigiai{
		Ta có $E = (6x^2 + 4x - 3x - 2) - (6x^2 + 6x) + 5x$ \\
$= 6x^2 + x - 2 - 6x^2 - 6x + 5x = -2$. \\
Vậy giá trị của biểu thức $E$ luôn bằng $-2$, không phụ thuộc vào biến $x$.
	}
\end{bt}

% nguồn: De_AI_Toan8_cuoi_hoc_ki_i_1 (nhập 21/11/2026)
\begin{bt}%[8D1H4-3]%[ma:f48f547e]
	Chứng minh rằng giá trị của biểu thức $F = (x - 2y)(x + 3y) - x(x + y) + 6y^2 + 10$ không phụ thuộc vào các biến $x$ và $y$.
	\loigiai{
		Ta có $F = (x^2 + 3xy - 2xy - 6y^2) - (x^2 + xy) + 6y^2 + 10$ \\
$= x^2 + xy - 6y^2 - x^2 - xy + 6y^2 + 10 = 10$. \\
Vậy giá trị của $F$ luôn bằng $10$, không phụ thuộc vào $x$ và $y$.
	}
\end{bt}

% nguồn: De_AI_Toan8_cuoi_hoc_ki_i_1 (nhập 21/11/2026)
\begin{bt}%[8D1H4-3]%[ma:a35b6213]
	Chứng minh rằng với $x \ne 0$, giá trị của biểu thức $G = (6x^3 - 3x^2) : (3x^2) - (4x^2 - 2x) : (2x)$ không phụ thuộc vào $x$.
	\loigiai{
		Với $x \ne 0$, ta có: \\
$(6x^3 - 3x^2) : (3x^2) = 2x - 1$; \\
$(4x^2 - 2x) : (2x) = 2x - 1$. \\
Suy ra $G = (2x - 1) - (2x - 1) = 0$. Vậy giá trị của $G$ không phụ thuộc vào $x$.
	}
\end{bt}

% nguồn: De_AI_Toan8_cuoi_hoc_ki_i_1 (nhập 21/11/2026)
\begin{bt}%[8D1H4-3]%[ma:8f4bc118]
	Chứng minh rằng giá trị của biểu thức $K = (x^2 - 2x + 3)(x - 4) - x(x^2 - 6x + 11)$ không phụ thuộc vào $x$.
	\loigiai{
		Ta có $(x^2 - 2x + 3)(x - 4) = x^3 - 4x^2 - 2x^2 + 8x + 3x - 12 = x^3 - 6x^2 + 11x - 12$. \\
Suy ra $K = (x^3 - 6x^2 + 11x - 12) - (x^3 - 6x^2 + 11x) = -12$. \\
Vậy giá trị của $K$ luôn bằng $-12$, không phụ thuộc vào $x$.
	}
\end{bt}

% nguồn: De_AI_Toan8_cuoi_hoc_ki_i_1 (nhập 21/11/2026)
\begin{bt}%[8D1H4-3]%[ma:bd66357e]
	Chứng minh rằng với mọi số nguyên $n$, giá trị của biểu thức $P = n(n + 5) - (n - 3)(n + 2)$ luôn chia hết cho $6$.
	\loigiai{
		Ta có $P = (n^2 + 5n) - (n^2 + 2n - 3n - 6) = n^2 + 5n - n^2 + n + 6 = 6n + 6 = 6(n + 1)$. \\
Vì $n \in \mathbb{Z}$ nên $n + 1 \in \mathbb{Z}$, suy ra $6(n + 1) \vdots 6$. Vậy $P$ luôn chia hết cho $6$ với mọi số nguyên $n$.
	}
\end{bt}

% nguồn: De_AI_Toan8_cuoi_hoc_ki_i_1 (nhập 21/11/2026)
\begin{bt}%[8D1H4-3]%[ma:c11f9aaf]
	Cho $x^2 + y^2 = 2$. Chứng minh rằng $2(x + 1)(y + 1) = (x + y)(x + y + 2)$.
	\loigiai{
		Biến đổi vế phải: $(x + y)(x + y + 2) = (x + y)^2 + 2(x + y) = x^2 + 2xy + y^2 + 2x + 2y$. \\
Thay $x^2 + y^2 = 2$ vào vế phải, ta được $2 + 2xy + 2x + 2y = 2(xy + x + y + 1)$. \\
Biến đổi vế trái: $2(x + 1)(y + 1) = 2(xy + x + y + 1)$. \\
Vậy vế trái bằng vế phải (điều phải chứng minh).
	}
\end{bt}

% nguồn: chuyen_dang (nhập 22/11/2026)
\begin{bt}%[8D1N4-3]%[ma:c7b7a3d7]
	Chứng minh rằng giá trị của biểu thức $K = (x + 5)^2 - x(x + 10)$ không phụ thuộc vào giá trị của biến $x$.
	\loigiai{
		Ta có $K = (x + 5)^2 - x(x + 10) = (x^2 + 10x + 25) - (x^2 + 10x) = x^2 + 10x + 25 - x^2 - 10x = 25$. \\
Vì $K = 25$ là một hằng số nên giá trị của $K$ không phụ thuộc vào giá trị của biến $x$.
	}
\end{bt}

% nguồn: chuyen_dang (nhập 22/11/2026)
\begin{bt}%[8D1H4-3]%[ma:8f4ca59e]
	Chứng minh giá trị của biểu thức sau không phụ thuộc vào giá trị của biến $x$:
	\[ A=(2x-1)^2-(2x+3)(2x-3)+4x. \]
	\loigiai{
		Ta có
		\[ A=\left(4x^2-4x+1\right)-\left(4x^2-9\right)+4x=4x^2-4x+1-4x^2+9+4x=10. \]
		Vậy $A=10$ với mọi $x$, tức là giá trị của $A$ không phụ thuộc vào $x$.
	}
\end{bt}

% nguồn: chuyen_dang (nhập 22/11/2026)
\begin{bt}%[8D1V4-3]%[ma:58ff5517]
% nguon: Archive/SP_KhoLatex/ID6_old/1K1-ot.tex #3782
%Câu 35
    Chứng minh biểu thức $\sin^2 x\cdot\tan x+4\sin^2 x-\tan^2 x+3\cos^2 x$ không phụ thuộc vào $x$
    \loigiai{
        \begin{eqnarray*}\sin^2 x\cdot\tan x+4\sin^2 x-\tan^2 x+3\cos^2 x&=&\left( \sin^2 x-1 \right)\tan^2 x+4\sin^2 x+3\cos^2 x\\
            &=&-\cos^2 x\cdot \tan^2 x+4\sin^2 x+3\cos^2x\\
            &=&-\sin^2 x+4\sin^2 x+3\left( 1-\sin^2 x \right)\\
            &=&3.\end{eqnarray*}}
\end{bt}

% nguồn: chuyen_dang (nhập 22/11/2026)
\begin{bt}%[8D1V4-3]%[ma:0074e448]
	
	Chứng minh rằng biểu thức sau đây không phụ thuộc vào $x$.
	\begin{enumEX}{2}
		\item  $A=\left( \sin x+\cos x \right)^2+\left( \sin x-\cos x \right)^2$.
		\item  $B=\sin^4 x-\cos^4 x-2\sin^2 x+1$.
	\end{enumEX}
	\loigiai{
		\begin{enumEX}{1}
			\item  $A=\left( \sin x+\cos x \right)^2+\left( \sin x-\cos x \right)^2=1+2\sin x\cos x+1-2\sin x\cos x=2$.\\
			\item  $\begin{aligned}[t]
				B&=\sin^4 x-\cos^4 x-2\sin^2 x+1=\left( \sin^2 x+\cos^2 x \right)\left( \sin^2 x-\cos^2 x \right)-2\sin^2 x+1\\
				&=1\left[ \sin^2 x-\left( 1-\sin^2 x \right) \right]-2\sin^2 x+1=0.
			\end{aligned}$
	\end{enumEX}}
\end{bt}

% nguồn: chuyen_dang (nhập 22/11/2026)
\begin{bt}%[8D1N4-3]%[ma:77dbe699]
	Chứng minh rằng biểu thức $P = x^2 - 6x + 14$ luôn nhận giá trị dương với mọi giá trị của $x$, từ đó tìm giá trị nhỏ nhất của $P$.
	\loigiai{
		Ta có $P = x^2 - 6x + 14 = (x^2 - 6x + 9) + 5 = (x - 3)^2 + 5$. \\
Vì $(x - 3)^2 \ge 0$ với mọi $x$, nên $P = (x - 3)^2 + 5 \ge 5 > 0$ với mọi $x$. \\
Dấu bằng xảy ra khi $x - 3 = 0 \Leftrightarrow x = 3$. \\
Vậy giá trị nhỏ nhất của $P$ là $5$, đạt được khi $x = 3$.
	}
\end{bt}
